% Outline: Discussion (1 page cap; MUST contain 6-month future work)
\section{Discussion}

Two configured quantities decide how much of the offline schema signal
survives: the decision timeout, which forfeits 34.6\% of decisions at the
documented 15\,ms and none at 75\,ms, and the running-batch cap, which
decides whether a queue forms at all. Uncapped, ordering changed nothing; capped at 16, the learned order
beats arrival order by 4.2\% while the free length heuristic runs 8.5\%
worse than doing nothing.

\subsection{Which Gateway Optimizations Apply, and Which We Did Not Test}
Three caches are in play. \emph{Semantic response caching}, the gateway's
main lever, was unavailable: its cache excludes streaming requests, and
every captured request streamed. \emph{Prefix KV caching} is a
different mechanism that E1's byte-identical schema prefix makes plausible:
with caching on and nothing else changed, pooled mean TTLT falls by 17\% on
the two stock-scheduler arms and 22--23\% on the four policy-base arms, at a
62\% peak hit rate. That run cached the whole stream, so it did not isolate
the schema's share. E3 reports the third, schema-embedding caching,
insufficient.

A six-stage run enabled the machinery the ordering rounds disabled: the
gateway's scheduler path, queue limits and scorer discipline, at 5.76
requests/s. Mean TTLT never left the run's drift band (largest consecutive
difference 70\,ms, 1.6\% of the 4.4\,s mean,
against an ABBA drift spread of up to 1.2\%), nothing was shed, and goodput
was flat. \emph{Queue-triggered scoring} consults the scorer only when the
gateway's admission queue forms; that queue did not form at this rate, and
the decision services
recorded zero scoring requests, so the run measures trigger-path
non-activation, not a free scorer.

\subsection{What the Negative Results Buy}
Identity features fail because tool combinations essentially never repeat in
training data; removing truncation produced no detected improvement; caching
fails because a $1.9\times$
speedup does not bridge a $48\times$ budget violation. Together: neither CPU rescue makes this BERT-base Ranker fit the
15\,ms default, and fitting it needs either the GPU placement
with a recalibrated budget (75\,ms covers every decision) or a path off the
critical path, which E6 makes more valuable still.

\subsection{The Baseline Was the Result}
The one large effect in our uncontended serving measurements did not come
from our contribution. Swapping the scheduler implementation while holding
arrival order fixed cut mean time-to-last-token by 44\%; changing the order
moved it by less than the session drifted over the same hours. An earlier
round, benchmarked against the stock scheduler alone, reported uniform gains
across every policy we tried. That should have raised suspicion: policies
that order differently should not improve identically.
Those gains came from the implementation underneath. The same mistake awaits
anyone who takes a serving engine's default scheduler as a baseline; the
control must share your implementation and differ only in the variable
claimed.

\subsection{Limitations}
Offline claims rest on one workload family (ToolACE) and one label model,
the live traces come from a single agent product, and the 75 captured
payloads calibrate the synthetic workload, so they cannot double as a
validation set. The replay inherits the capture's compressed job-size
variance, which may leave ordering little room to separate.

One choice bounded the first three ordering rounds: vLLM's running-batch cap
sat at its 256-slot default. By Little's law our own
numbers put mean occupancy at roughly eleven requests in flight at
3.96\,requests/s and nineteen at 6.4, under a thirteenth of that cap. Mean
occupancy does not forbid queueing, and our logs show it did not: bursts put
11.9--12.5\% of requests behind another, but those queues drained within a
scheduling step instead of accumulating. Our calibration targeted completion
(a 95\% success criterion), not contention. E6 removes that bound and the
ordering result flips, so every ordering number here is conditional on its
regime. We did not vary how contention arises: E6 constrains the batch, while contexts large enough to
exhaust the KV cache would reach it another way.

The main matrix rests on three replay repetitions inside a single engine launch
per arm, and E6 on one cold launch per arm. The Ranker predicts decode
length
only, so prefill cost is unmodeled; the offline split is row-level; the
replay is open-loop; S2 is underpowered at $n{=}78$; and the gate's
confidence is a measured $\tau$ floor, not a calibrated probability. Serving
hardware is a vendor-modified 48\,GB RTX~4090.

\subsection{Future Work}
Six directions would occupy the next six months, in the order we would take
them. First, add an oracle-order arm sorting by measured length, separating
residual ranking error from the mechanism E6 exposes, and test whether the
gate's cost under contention can be recovered by letting abstained requests
keep a bounded slot rather than a fixed one. Second, validate
queue-triggered scoring under a load that forms the gateway queue. Third,
close the predictor-cost gap by distillation. Fourth, grow the S2 stratum to
statistical power. Fifth, move to closed-loop session benchmarks. Sixth, explore continual fine-tuning on production traffic under explicit
leakage discipline: training and evaluation traffic must never mix.
